Poniższa dokumentacja opisuje interfejsy API oraz komponenty zaimplementowane w kodzie źródłowym systemu, zweryfikowane na dzień przygotowania niniejszego załącznika.
\subsection*{Moduły systemu}

System został zbudowany z modułów, które wykorzystują model AI do świadczenia usług komercyjnych. System składa się z następujących modułów:

\begin{itemize}
\item \textbf{IQText} -- moduł automatycznie wykrywający i klasyfikujący dane, umożliwiający identyfikację typów informacji na podstawie ich treści i kontekstu. Analizuje treść dokumentów wejściowych, dzieląc ją na fragmenty z uwzględnieniem struktury tekstu prawnego (numeracja, paragrafy, artykuły, rozdziały), a następnie przy użyciu modelu językowego (LLM) ekstrahuje z niej zdefiniowane pola danych. Wyniki cząstkowe są scalane w jeden spójny obiekt danych, walidowane względem schematu XSD i udostępniane w formacie XML/JSON poprzez REST API, wraz z interfejsem umożliwiającym użytkownikowi weryfikację i edycję wyekstrahowanych danych (szczegóły w podrozdziałach 1.1 i 1.2).
\item \textbf{Anonymizer} -- moduł anonimizacji danych w dokumentach, zastępujący dane poufne tokenami w celu ochrony informacji wrażliwych. Automatycznie wykrywa fragmenty danych wrażliwych, łącząc podejście regułowe dla danych o ściśle zdefiniowanym formacie (np. PESEL, NIP, IBAN) z modelem klasyfikującym typu NER/LLM dla danych o zmiennym formacie (imiona, nazwiska, adresy), a następnie ukrywa je jedną z metod dobieranych w zależności od wymaganej odwracalności: nieodwracalnym maskowaniem/redakcją, odwracalną tokenizacją z osobno przechowywanym mapowaniem, jednokierunkowym haszowaniem z solą lub szyfrowaniem AES-256 z kluczem zarządzanym przez dedykowany mechanizm KMS/HSM (szczegóły w podrozdziałach 2.1 i 2.2).
\item \textbf{DLP} -- rozproszony system Atende DLP (\emph{Data Loss Prevention}) wykrywający i przeciwdziałający nieautoryzowanemu wyciekowi danych wrażliwych ze stacji roboczych. Składa się z trzech niezależnie rozwijanych modułów: \emph{DLP Agent} (usługa Windows w C\#/.NET, monitorująca operacje plikowe mechanizmem ETW oraz zawartość schowka systemowego), \emph{DLP Backend Server} (serwer Java/Spring Boot z bazą MySQL/MariaDB, jedyny moduł wystawiający publiczne REST API panelu administracyjnego oraz API przyjmujące dane od agentów, pracujący w profilu \emph{collector} -- u klienta -- lub \emph{aggregator} -- centralnie agregującym alerty z wielu instalacji) oraz \emph{DLP Parser} (skrypty wsadowe w Pythonie, uruchamiane cyklicznie przez backend, ekstrahujące treść plików, w tym poprzez OCR, i wykrywające w niej dane wrażliwe regułami oraz wtyczkami). Wykryte naruszenia rejestrowane są jako alerty w panelu administracyjnym, a alerty krytyczne dodatkowo rozsyłane e-mailem (szczegóły w podrozdziale 3).
\item \textbf{Etykietowanie} -- moduł wspierający prace treningowe modelu poprzez oznaczanie i przygotowywanie danych do procesu uczenia. Udostępnia narzędzie do kontekstowego, ręcznego klasyfikowania dokumentów przez zespół analityczny oraz zarządzania zbiorami etykiet -- zarówno kategorii dokumentów, jak i typów danych wrażliwych -- a także przygotowywania mapowań typu NER (Named Entity Recognition), łączących typy danych, kontekst ich występowania oraz przypisane etykiety. Tak przygotowane zbiory stanowią materiał wejściowy dla modeli klasyfikujących i ekstrahujących dane wrażliwe wykorzystywanych w pozostałych modułach systemu (szczegóły w podrozdziale 4).
\end{itemize}

W poniższym rozdziale opisano kolejno każdy z powyższych modułów, a na końcu -- wyniki testów prototypu wypracowanego w toku ich integracji (zadania R1--R3).

\subsection*{Moduł 1: IQText}
\addcontentsline{toc}{subsection}{Moduł 1: IQText}

\subsubsection*{1.1. Backend systemu IQText}

Backend zaimplementowano w technologii \textbf{Flask} (Python), z wykorzystaniem bibliotek \texttt{flask-cors} (obsługa CORS), \texttt{flask\_apscheduler} (zadania cykliczne) oraz \texttt{flasgger} (automatyczna dokumentacja Swagger/OpenAPI dostępna pod ścieżką \texttt{/apidocs} działającej instancji serwera). Aplikacja przechowuje dane w bazie SQLite oraz na systemie plików (katalogi dokumentów wejściowych, sparsowanych i wynikowych, konfigurowane zmiennymi środowiskowymi).

\begingroup
\small
\begin{longtable}{p{0.09\textwidth} p{0.20\textwidth} p{0.30\textwidth} p{0.32\textwidth}}
\caption{Endpointy REST API backendu IQText} \\
\toprule
\rowcolor{tablehead!15}
Metoda & Ścieżka & Parametry & Odpowiedź \\
\midrule
\endfirsthead
\toprule
\rowcolor{tablehead!15}
Metoda & Ścieżka & Parametry & Odpowiedź \\
\midrule
\endhead
GET & \texttt{/api} & brak & Prosty komunikat tekstowy potwierdzający dostępność usługi (adres IP klienta). \\
GET & \texttt{/api/get-xml} & query: \texttt{project} (numer projektu) & Treść XML wygenerowana dla wskazanego projektu lub komunikat błędu. \\
GET & \texttt{/api/get-json} & query: \texttt{project} & JSON: dane wyekstrahowane przez model AI (\texttt{data}), dane po edycji użytkownika (\texttt{editedData}), znaczniki błędów i czasu edycji. \\
POST & \texttt{/api/save-json} & body JSON: \texttt{json}, \texttt{project\_number} & JSON z potwierdzeniem zapisu edytowanych danych lub komunikatem błędu. \\
POST & \texttt{/api/login} & body JSON: \texttt{email}, \texttt{password} & JSON z potwierdzeniem/odrzuceniem logowania (hasło weryfikowane wobec skrótu bcrypt). \\
GET & \texttt{/api/projects} & brak & JSON z listą numerów projektów zarejestrowanych w systemie. \\
POST & \texttt{/api/add-\allowbreak document} & form-data: \texttt{project}, \texttt{file} (jeden lub więcej plików) & JSON z informacją o zapisanych plikach lub komunikatem błędu. \\
POST & \texttt{/api/szop/\allowbreak analyze} & body JSON: \texttt{request\_id} & JSON potwierdzający przyjęcie żądania do przetworzenia (integracja z usługą zewnętrzną SZOP) lub komunikat błędu. \\
POST & \texttt{/api/generate\_\allowbreak xml} & body JSON: \texttt{content} & Dokument XML zweryfikowany wobec schematu XSD lub JSON z opisem błędu walidacji. \\
POST & \texttt{/api/get\_xml\_\allowbreak szop} & query: \texttt{project} & Odpowiedź przekazana bez zmian z usługi zewnętrznej SZOP (XML lub JSON). \\
POST & \texttt{/api/init-xml} & query: \texttt{project} & JSON z informacją o zapisanych plikach inicjujących projekt. \\
\bottomrule
\end{longtable}
\endgroup

Niezależnie od endpointów HTTP, backend uruchamia zadanie cykliczne (harmonogram co 10 sekund) realizujące pełny potok przetwarzania nowych dokumentów: parsowanie treści plików, wywołanie modelu językowego (LLM) do ekstrakcji danych, wygenerowanie i walidację pliku XML względem schematu XSD, a następnie przekazanie wyniku do usługi zewnętrznej.



\subsubsection*{1.2. Moduł ekstrakcji danych z dokumentów (\texttt{contract\_extractor} / EL-GENERATOR / EM-GENERATOR)}

Moduł nie udostępnia własnego, oddzielnego API REST (poza wyjątkiem opisanym w podrozdziale 1.3) -- jest wywoływany bezpośrednio przez backend IQText w ramach zadania cyklicznego. Jego działanie obejmuje:

\begin{itemize}
    \item wygenerowanie zadań/promptów dla modelu LLM na podstawie struktury schematu XSD (klasa \texttt{XSDTaskGenerator});
    \item podział treści dokumentu na fragmenty (\emph{chunking}) z uwzględnieniem struktury prawnej tekstu (numeracja, paragrafy, artykuły, rozdziały);
    \item wykrycie obszarów danych występujących w dokumencie (DORA / RODO / NDA / NIS-UKSC / INNE);
    \item właściwą ekstrakcję pól danych dla każdego fragmentu przy użyciu modelu LLM (lokalny model uruchamiany przez Ollama, domyślnie \texttt{llama3.3:70b}, lub API OpenAI);
    \item scalenie wyników z wielu fragmentów dokumentu w jeden spójny obiekt JSON, z wykorzystaniem LLM do wyboru najlepszej z kilku kandydujących odpowiedzi;
    \item filtrowanie z wynikowej listy kanałów notyfikacji wpisów odpowiadających adresom e-mail w domenie firmowej oraz osobom z opcjonalnej listy wykluczeń (plik CSV) -- jest to reguła biznesowa ograniczająca sugerowanie pracowników własnej organizacji jako kontaktów klienta, nie mechanizm anonimizacji danych osobowych klienta.
\end{itemize}


\subsubsection*{1.3. Serwer pomocniczy modułu EM-GENERATOR}

Moduł EM-GENERATOR udostępnia dodatkowo lekki serwer HTTP (Flask), pośredniczący w komunikacji z lokalnym modelem LLM:

\begin{longtable}{p{0.09\textwidth} p{0.20\textwidth} p{0.30\textwidth} p{0.32\textwidth}}
\toprule
\rowcolor{tablehead!15}
Metoda & Ścieżka & Parametry & Odpowiedź \\
\midrule
\endhead
POST & \texttt{/search} & body JSON zawierający treść zapytania/promptu oraz fragment analizowanego tekstu & Surowa odpowiedź wygenerowana przez lokalny model LLM (Ollama, domyślnie \texttt{llama3:latest}). \\
\bottomrule
\end{longtable}

Serwer ten nie posiada mechanizmu uwierzytelniania, limitów zapytań ani walidacji pochodzenia żądania -- jest przeznaczony do użytku lokalnego/deweloperskiego.

\subsection*{Moduł 2: Anonymizer}
\addcontentsline{toc}{subsection}{Moduł 2: Anonymizer}

\subsubsection*{2.1. Narzędzie redakcji dokumentów (Anonymizer)}

Moduł Anonymizer to samodzielne narzędzie zaimplementowane w technologii \textbf{FastAPI}, wykorzystujące bibliotekę \textbf{PyMuPDF} (\texttt{fitz}) do operacji na warstwie tekstowej plików PDF.

\begin{longtable}{p{0.07\textwidth} p{0.14\textwidth} p{0.35\textwidth} p{0.35\textwidth}}
\caption{Endpointy narzędzia Anonymizer} \\
\toprule
\rowcolor{tablehead!15}
Metoda & Ścieżka & Parametry & Odpowiedź \\
\midrule
\endfirsthead
\toprule
\rowcolor{tablehead!15}
Metoda & Ścieżka & Parametry & Odpowiedź \\
\midrule
\endhead
GET & \texttt{/} & brak & Strona HTML z interfejsem umożliwiającym wskazanie par ,,znajdź / zamień'' dla wczytanego dokumentu. \\
GET & \texttt{/pdf} & brak & Zawartość aktualnie wczytanego dokumentu PDF (plik demonstracyjny). \\
POST & \texttt{/process} & form-data: listy \texttt{find\_text} oraz \texttt{replace\_text} & Zmodyfikowany plik PDF, w którym wskazane fragmenty tekstu zostały fizycznie zredagowane (usunięte z warstwy treści) i zastąpione podanym tekstem zastępczym. \\
\bottomrule
\end{longtable}

\textbf{Zasada działania:} interfejs użytkownika pozwala wskazać dokładne fragmenty tekstu do ukrycia lub zezwolić mechanizmom AI na ich wskazanie; dla każdego wskazanego fragmentu system generuje losowy token zastępczy o długości oryginalnego tekstu. Backend wyszukuje dokładne dopasowania wskazanego tekstu na każdej stronie dokumentu, stosuje redakcję (fizyczne usunięcie z warstwy treści PDF), a następnie wstawia w to samo miejsce tekst zastępczy.

\textbf{Zakres funkcjonalny obecnej wersji narzędzia:}
\begin{itemize}
    \item  automatyczna detekcjia danych wrażliwych -- użytkownik samodzielnie wskazuje fragmenty tekstu do zamiany lub pozwala to na wykonanie automatowi;
    
    \item  mechanizmy szyfrowania danych w dokumencie PDF, tokenizacja odwracalna lub haszowania -- opcje te są zaznaczone w interfejsie użytkownika jako punkt wyjścia do procesu anonimizacji.
\end{itemize}


\subsubsection*{2.2. Specyfikacja funkcjonalna modułu anonimizacji}

Moduł anonimizacji automatycznie wykrywa dane wrażliwe w dokumentach przetwarzanych przez system i ukrywa je jedną z dostępnych metod (maskowanie, tokenizacja, haszowanie lub szyfrowanie AES), zależnie od tego, czy dane mają pozostać nieodwracalne, czy odtwarzalne przez uprawnionego użytkownika. 

\subhead{Cel modułu}

Celem modułu anonimizacji jest automatyczne wykrywanie danych wrażliwych w dokumentach przekazywanych do systemu oraz ich nieodwracalne lub odwracalne ukrycie, zgodnie z zasadą \emph{privacy by design/by default} (art. 25 RODO), przed dalszym przetwarzaniem, prezentacją lub archiwizacją dokumentu.

\subhead{Zakres wykrywanych kategorii danych}

Moduł wykorzystuje istniejące zbiory etykiet danych wrażliwych opracowane w zadaniu B3 (33 etykiety) jako punkt wyjścia, uzupełnione o wzorce charakterystyczne dla polskiego i europejskiego kontekstu prawnego, w tym co najmniej:

\begin{itemize}
    \item dane identyfikacyjne osoby: imię, nazwisko, numer PESEL, numer dowodu osobistego/paszportu;
    \item dane kontaktowe: adres zamieszkania/korespondencyjny, adres e-mail, numer telefonu;
    \item dane identyfikacyjne podmiotów: NIP, REGON, numer KRS, numer rachunku bankowego (IBAN);
    \item daty istotne z punktu widzenia ochrony danych (data urodzenia) oraz retencji (rozpoznawane przez model opracowany w zadaniu B5);
    \item dane finansowe (kwoty, waluty) rozpoznawane regułowo, analogicznie do istniejącej walidacji kwot w module ekstrakcji danych.
\end{itemize}

Detekcja łączy podejście regułowe (wyrażenia regularne dla danych o ściśle zdefiniowanym formacie, np. PESEL, NIP, IBAN) z modelem klasyfikującym typu NER/LLM, wykorzystującym mapowania i etykiety wypracowane w zadaniu B3, w celu rozpoznawania danych o zmiennym formacie (imiona, nazwiska, adresy) w kontekście dokumentu.

\subhead{Architektura modułu}

\begin{itemize}
    \item \textbf{Warstwa detekcji} -- komponent analizujący treść dokumentu i zwracający listę wykrytych fragmentów wraz z: kategorią danych, lokalizacją w dokumencie (strona, współrzędne/zakres znaków) oraz poziomem pewności detekcji.
    \item \textbf{Warstwa anonimizacji} -- komponent stosujący do każdego wykrytego fragmentu jedną z metod ukrycia danych, wybieraną w zależności od dalszego przeznaczenia danych:
    \begin{itemize}
        \item \emph{Maskowanie/redakcja} -- zastąpienie wartości stałym symbolem (np. \texttt{[DANE\_USUNIĘTE]}) lub gwiazdkami; nieodwracalne, najniższy koszt implementacji; rozszerza już zrealizowany mechanizm redakcji PDF o automatyczne źródło fragmentów do ukrycia (z warstwy detekcji, zamiast ręcznego wskazywania przez użytkownika).
        \item \emph{Tokenizacja odwracalna} -- zastąpienie oryginalnej wartości losowym identyfikatorem (tokenem), przy czym mapowanie token~$\rightarrow$~wartość oryginalna przechowywane jest osobno, w zaszyfrowanej postaci, z dostępem ograniczonym do uprawnionych ról.
        \item \emph{Haszowanie jednokierunkowe} -- dla przypadków, w których wymagana jest wyłącznie możliwość porównania dwóch wartości bez ich odtwarzania (np. deduplikacja rekordów), z wykorzystaniem funkcji skrótu z solą (SHA-256 + losowa sól per rekord), analogicznie do stosowanego w systemie mechanizmu haszowania haseł (bcrypt).
        \item \emph{Szyfrowanie odwracalne (AES)} -- dla przypadków wymagających pełnego odtworzenia oryginalnej wartości przez uprawnionego użytkownika: szyfrowanie symetryczne AES-256 w trybie uwierzytelnionym (AES-256-GCM), z kluczem zarządzanym przez dedykowany mechanizm zarządzania kluczami (zewnętrzny KMS/HSM), nieprzechowywanym w tej samej bazie co dane zaszyfrowane.
    \end{itemize}
    \item \textbf{Warstwa zarządzania kluczami i mapowaniami} -- osobny magazyn (odizolowany od bazy danych `Metadata`/`Results`/`Json`) przechowujący klucze szyfrujące oraz mapowania tokenizacyjne, z rejestrowaniem (audytem) każdego dostępu do funkcji odtworzenia danych oryginalnych.
    \item \textbf{Warstwa API} -- zestaw endpointów udostępniających funkcjonalność modułu pozostałym komponentom systemu (backend IQText, wizualizator IQText-GUI).
\end{itemize}

\subhead{API modułu}

Poniższe endpointy stanowią interfejs programistyczny modułu, zaimplementowany spójnie ze stylem API backendu IQText (Flask, prefiks \texttt{/api}).

\begingroup
\small
\begin{longtable}{p{0.06\textwidth} p{0.27\textwidth} p{0.29\textwidth} p{0.29\textwidth}}
\caption{Endpointy modułu anonimizacji} \\
\toprule
\rowcolor{tablehead!15}
Metoda & Ścieżka & Parametry wejściowe & Odpowiedź \\
\midrule
\endfirsthead
\toprule
\rowcolor{tablehead!15}
Metoda & Ścieżka & Parametry wejściowe & Odpowiedź \\
\midrule
\endhead
POST & \texttt{/api/anonymize/\allowbreak detect} & identyfikator dokumentu lub treść dokumentu & lista wykrytych fragmentów danych wrażliwych wraz z kategorią, lokalizacją i poziomem pewności \\
POST & \texttt{/api/anonymize/\allowbreak apply} & identyfikator dokumentu, lista fragmentów do ukrycia, wybrana metoda (maskowanie / tokenizacja / haszowanie / AES) & zanonimizowany dokument lub jego identyfikator; w przypadku metod odwracalnych -- identyfikator zapisanego mapowania/klucza \\
POST & \texttt{/api/anonymize/\allowbreak reveal} & identyfikator zanonimizowanego fragmentu/dokumentu, autoryzacja uprawnionego użytkownika & oryginalna wartość danych (wyłącznie dla metod odwracalnych, z zapisem w dzienniku audytowym) \\
GET & \texttt{/api/anonymize/\allowbreak methods} & brak & lista dostępnych metod anonimizacji wraz z opisem ich odwracalności \\
\bottomrule
\end{longtable}
\endgroup
